\section{The remaining scale-dependent comparison}
\label{ctn-next}

For the bare composite, the new identification reduces
the leading coefficient, up to the uniformly bounded
coarse effective-Hessian term, to the constant-curvature
jets of
\[
 \log\det(D_b^TD_b)
 -\frac12\log\det
       \bigl(S_{\mathrm W}''+D_bD_b^T+aC_K^TC_K\bigr).
\]
This is a statement about the actual diagonal trace density
and its Maxwell moment. The formal Cartan probe makes
the constrained-current curvature vanish in the required
jets. It does not make either determinant individually
uniformly massive at the final scale.

The next estimate must compare this combined expression
to a fully specified background-Feynman vector and ghost
reference with infrared scale \(K^{-1}\), preserving its
coarse-coordinate, color and action normalization. In
particular, changing the based scalar domain, the vector
constraint completion and the infrared regulator must
be controlled together. The based-ghost Rayleigh bound
from Notebook 32 still rules out a separate inverse
bound of order \(K^2\).

A sufficient result would give a constant \(C\), independent
of \(K\), bounding the difference of the extracted
Maxwell coefficients. The specified continuum reference
now has an explicit positive logarithmic law and bounded
finite part. The exact actual heat representation identifies
the remaining comparison integral; the three interval
estimates in Lemma~\ref{ahc-band-lemma} would bound it.
Those estimates do not follow from the fixed-\(K\)
convergence used here. The actual based-ghost part of
the short-time interval is now bounded uniformly by
the closed-walk argument; its vector counterpart and
the other intervals remain. Uniformity over
the structural systems needed by the finite drift-transfer
theorem must also be retained.

Once positive averaged drift is established with that
scope, the conditional finite shooting theorem applies.
The full four-dimensional existence objective still
requires the other field sectors, the exact global
measure representation, sourced continuum bounds,
nontriviality, Euclidean symmetry, positivity and
reconstruction.

\begin{thebibliography}{9}
\bibitem{vassilevichctn}
D.~V.~Vassilevich,
\emph{Heat kernel expansion: user's manual},
Physics Reports 388 (2003), 279--360,
\href{https://arxiv.org/abs/hep-th/0306138}{arXiv:hep-th/0306138}.
Sections 3.4, 4.2.1 and 8.3 specify the background vector
and ghost operators, their combined heat coefficient,
and the constant magnetic scalar heat kernel.
This source is not used as a proof of the actual constrained
lattice comparison.
\bibitem{priorctn}
\emph{Toward a four-dimensional Yang--Mills construction},
research Notebooks 15--32 in this workspace.
The actual local minimum, strict mixed kernels, classical
envelope, constrained Gaussian, gauge identities, determinant
bridge and Maxwell normalization are the inputs used here.
\bibitem{sourcectn}
E.~Sussman and collaborators,
\href{https://github.com/ethanwsussman/Yang--Mills-LLM-experiment}
{Yang--Mills-LLM-experiment}, pinned deposited accounts at
commit \texttt{e392ef45013a581664bab76b90867170b30fe10e}.
The centered ordinary average and its selected
endpoint-equivariant mean are kept unchanged.
\end{thebibliography}
\end{document}
